\section{Chapter~\ref{sec:nora}. \nt{NORA}}

The organization of the \nt{NORA} system, its two modes, its link to the pupil (not only through \nt{NORA}), the rise in signal-to-background ratio and the inverted U in behavior are measured directly; the multiplicative gain $g$ is a model; that gain switches the choice mode and acts as an inverse temperature are conclusions of the chapter; the benefit of adjusting the gain is computed with the book's model; ``\nt{NORA} as the exploration knob'' and ``\nt{NORA} as an interrupt'' are hypotheses.

{\footnotesize
\setlength{\tabcolsep}{3pt}\renewcommand{\arraystretch}{1.15}
\begin{longtable}{>{\raggedright}p{0.5cm}>{\raggedright}p{4.0cm}>{\raggedright}p{5.6cm}>{\raggedright}p{2.3cm}>{\raggedright\arraybackslash}p{1.7cm}}
\toprule
No. & Claim & Experiment and what was measured & Source & Status \\
\midrule\endfirsthead
\toprule
No. & Claim & Experiment and what was measured & Source & Status \\
\midrule\endhead
1 & A human has a few tens of thousands of \nt{NORA} neurons, almost all in one group & Serial sections of the human brainstem stained for tyrosine hydroxylase: $53\,900$ cells in the locus coeruleus and another $6\,260$ in the adjacent subnucleus. & \cite{Baker1989LC} & measured \\
2 & The outputs spread through almost the whole brain, but parts of the group partly serve different regions & Viral-genetic labeling in mice: \nt{NORA} neurons of the locus coeruleus receive inputs from wide areas of the brain and send outputs to wide areas of the brain and spinal cord. In rats: a group projecting to the amygdala enhances fear learning, and a separate group projecting to the prefrontal cortex enhances its extinction. & \cite{Schwarz2015LC, Uematsu2017LC, Poe2020} & measured \\
3 & Tonic mode: steady spikes at rates from a fraction of a spike to a few per second, changing with wakefulness & Rats, recordings of locus coeruleus neurons across the sleep--wake cycle: the rate is highest in waking, lower in slow-wave sleep, and almost zero in REM sleep; rate changes precede state transitions. & \cite{AstonJonesBloom1981} & measured \\
4 & Phasic mode: a short burst to a task-relevant event, at roughly the moment of decision & Monkeys, a task of detecting a rare target: all locus coeruleus neurons respond with a burst only to the target, $\sim90\ms$ after it and $\sim200\ms$ before the hand response; bursts are weaker when performance is poor. In a two-choice task the burst is more tightly locked to the response than to the stimulus, and it occurs on both correct and error trials. & \cite{AstonJones1994vigilance, Clayton2004LC} & measured \\
5 & Pupil diameter is an imprecise test point for \nt{NORA}: it follows \nt{NORA}, \nt{ACET} and other regions & Monkeys: locus coeruleus spikes, down to single spikes of one cell, predict pupil dilation; a similar but later relation exists in the colliculi and the cingulate cortex. Mice: pupil fluctuations follow the activity of both noradrenergic and cholinergic outputs in the cortex. & \cite{Joshi2016, Reimer2016} & measured \\
6 & Noradrenaline suppresses background activity more than evoked activity; the multiplicative gain~\eqref{eq:gain} is a model that captures this effect & Awake monkeys, auditory cortex, iontophoresis of noradrenaline: background activity of neurons is suppressed more than the response to conspecific calls, so the signal-to-noise ratio rises. The multiplicative sigmoid~\eqref{eq:gain} is taken from a network model; a literal multiplication of the slope has not been measured. & \cite{Foote1975NE, ServanSchreiber1990} & measured; $g$ is a model \\
7 & Gain raises $d'$ only against noise after the amplifier~\eqref{eq:gainsnr} (Proposition~\ref{prop:gainnoise}) & Derived in the chapter; in the network model, gain improves signal detection. There is no experiment that separates noise before and after the amplifier and tests this distinction. & derived in the chapter; \cite{ServanSchreiber1990} & theory \\
8 & The boundary $g\beta=1$ switches the choice network from ``deliberating'' to ``decided''~\eqref{eq:wtagain} (Proposition~\ref{prop:gainwta}, Fig.~\ref{fig:pitchfork}) & Linear analysis of two mutually inhibiting groups; derived in the chapter. There are no direct measurements of this threshold in the brain. & derived in the chapter & theory \\
9 & The gain is the inverse temperature of choice~\eqref{eq:softmax} & With logistic noise after the amplifier, the choice probability is softmax with $T=\sigma/g$; derived in the chapter. & derived in the chapter & theory \\
10 & \nt{NORA} sets how much to explore: high tonic activity means small effective $g$ (exploration), a phasic burst at the moment of decision means large $g$ (decision) & Humans: the baseline pupil is wider before a random choice than before choosing the best known option; people with wider pupils are more inclined to explore. Rats: switching into a random-choice mode is controlled by \nt{NORA} input to the anterior cingulate cortex. That the burst raises the effective $g$ has not been measured directly. & \cite{JepmaNieuwenhuis2011, Tervo2014stochastic, AstonJones2005} & hypothesis \\
11 & Reward depends on $g$ as an inverted U; the best $g$ is smaller in a fast-changing world (Proposition~\ref{prop:explore}, Fig.~\ref{fig:invertedU}) & \texttt{models/nora\_explore.py}, $60$ runs of $4000$ trials. A change once per $1000$ trials: best $g=16$, fraction $0.976$; once per $20$: $g=8$, fraction $0.886$; at $g=0.5$, $0.77$. The rerun matches the chapter. & \texttt{models/\allowbreak nora\_\allowbreak explore.py} & model \\
12 & The inverted U in behavior: best performance at moderate tonic activity with sharp bursts & Monkeys, the same task as in row~4: during periods of good performance the tonic rate is moderate and bursts to the target are sharp; at a high tonic rate there are no bursts and more false responses. Changes in tonic and evoked activity follow fluctuations in performance. & \cite{Usher1999LC, AstonJones2005} & measured \\
13 & \nt{NORA} reports unexpected uncertainty, \nt{ACET} expected uncertainty & A theory of Bayesian inference with two kinds of uncertainty; the cited papers contain no direct experiment separating these signals by transmitter. & \cite{YuDayan2005} & hypothesis \\
14 & After an abrupt change people learn faster, and the pupil dilates & Humans, a prediction task with sudden changes: brief pupil changes track the estimated probability of a change and, together with the baseline pupil, predict how strongly new data change the estimate (the learning rate); pupil changes unrelated to light and task alter this rate. That the pupil here reflects \nt{NORA} specifically is an inference from the indirect data of row~5. & \cite{Nassar2012} & measured \\
15 & The phasic burst works as an interrupt: the network resets its state & There is no direct experiment in which a \nt{NORA} burst resets the network's state; only bursts to important events have been measured (row~4). & \cite{BouretSara2005} & hypothesis \\
16 & Surprise-driven adjustment of $g$ or of the learning rate is hardly better than the best constant settings & \texttt{models/nora\_explore.py}, a world with epochs of $1000$ trials (a change once per $1000$ and once per $20$). Best constant $g=8$: $0.925$; adjusting $g$: up to $0.927$; adjusting the learning rate: up to $0.926$; constant learning rate $0.4$: $0.928$. The rerun matches the chapter. & \texttt{models/\allowbreak nora\_\allowbreak explore.py} & model \\
\bottomrule
\end{longtable}}
